\chapter{Typing and definitional equality}
\label{chap:typing}

\section*{At a glance}
\begin{itemize}
\item \lead{Goal} State Lean4Lean's typing and definitional equality judgment and prove its
  structural metatheory.
\item \lead{Main results} Weakening (\ref{thm:isdefeq-weakn}); Substitution
  (\ref{thm:isdefeq-instn}); Regularity of types (\ref{thm:isdefeq-istype}); a well-formed
  environment is ordered (\ref{thm:wf-ordered}) and strongly typed (\ref{thm:wf-strong},
  \ref{thm:wf-orderedstrong}).
\item \lead{Status} \stProved{} the judgment and its structural metatheory; \stSorry{}
  \ref{thm:wf-patsstrong} (subject reduction for $\iota$) and \ref{thm:injectivity-open}
  (three inversion principles, pre-existing); \stTainted{} \ref{thm:wf-orderedstrong} and
  \ref{thm:foralle-inv-derived}.
\item \lead{Contribution} iota: the \texttt{pat} rule and admissibility
  (\ref{rule:isdefeq-pat}--\ref{def:patwf}); a \texttt{pat} case in twelve inductions; the
  well-formed-to-strong bootstrap (\ref{def:patsstrong}--\ref{thm:wf-orderedstrong}), open at
  \ref{thm:wf-patsstrong}. trproj: none.
\item \lead{Lean files} \texttt{Lean4Lean/Theory/Typing/Basic.lean}, \texttt{Env.lean},
  \texttt{Lemmas.lean}, \texttt{Meta.lean}, \texttt{QuotLemmas.lean}, \texttt{EnvLemmas.lean},
  \texttt{Injectivity.lean}.
\item \lead{Thesis} \texttt{axioms.tex} \S2.1--\S2.2, \S2.5--\S2.7; \texttt{typesys.tex}
  Weakening, Substitution, Regularity; \texttt{unique.tex} Regularity of reductions and
  Definitional inversion.
\end{itemize}

\section{The judgment}
\label{sec:typing-basic}

\begin{itemize}
\item $E$: environment; $U$: universe arity; $\Gamma$: de Bruijn context.
\item $\Gamma \vdash e_1 \equiv e_2 : A$ is \texttt{VEnv.IsDefEq} $E\,U\,\Gamma\,e_1\,e_2\,A$;
  $\Gamma \vdash e : A$ its diagonal (\texttt{HasType}).
\item $E \le E'$: the extension order of Chapter~\ref{chap:syntax}.
\end{itemize}

\begin{definition}[de Bruijn variable lookup]
\label{def:lookup}
\lean{Lean4Lean.Lookup}
\leanok
\uses{def:liftn}
\srcloc{Lean4Lean/Theory/Typing/Basic.lean}{6}
\thesisref{axioms.tex \S2.1, the variable rule}

\lead{In short} \texttt{Lookup}\,$\Gamma\,i\,A$: the \hl{$i$-th entry of $\Gamma$}, lifted
over the $i$ binders crossed, is $A$.
\begin{itemize}
\item \texttt{zero}: the head of $A :: \Gamma$ is $A$ lifted once; \texttt{succ}: recurse.
\end{itemize}
\end{definition}

\begin{definition}[Typing and definitional equality]
\label{def:isdefeq}
\contrib{mixed}
\lean{Lean4Lean.VEnv.IsDefEq}
\leanok
\uses{def:lookup, struct:venv, ind:vexpr, def:vlevel-wf, inductive:pattern-matches, def:pattern-check-realizes, inductive:pattern-rhs}
\srcloc{Lean4Lean/Theory/Typing/Basic.lean}{18}
\thesisref{axioms.tex \S2.1 and \S2.2 merged; \S2.6.4 for \texttt{pat}}

\lead{In short} One relation $\Gamma \vdash e_1 \equiv e_2 : A$, merging \hl{typing and
definitional equality}; every row below is such a judgment in the ambient $\Gamma$.
\begin{center}
\begin{tabular}{lp{0.45\linewidth}p{0.35\linewidth}}
\toprule
Rule & Premises & Conclusion \\
\midrule
\texttt{bvar} & \texttt{Lookup}\,$\Gamma\,i\,A$ & $\mathrm{bvar}\,i\equiv\mathrm{bvar}\,i:A$ \\
\texttt{symm} & $e\equiv e':A$ & $e'\equiv e:A$ \\
\texttt{trans} & $e_1\equiv e_2:A$, $e_2\equiv e_3:A$ & $e_1\equiv e_3:A$ \\
\texttt{sortDF} & $l,l'$ level-wf, $l\approx l'$ & $\mathrm{Sort}\,l\equiv\mathrm{Sort}\,l':\mathrm{Sort}(l{+}1)$ \\
\texttt{constDF} & $c$ registered, $ls,ls'$ level-wf, pointwise $\approx$ & $\mathrm{const}\,c\,ls\equiv\mathrm{const}\,c\,ls':\tau_c[ls]$ \\
\texttt{appDF} & $f\equiv f':\forall A.B$, $a\equiv a':A$ & $f\,a\equiv f'\,a':B[a]$ \\
\texttt{lamDF} & $A\equiv A':\mathrm{Sort}\,u$, $A{::}\Gamma\vdash b\equiv b':B$ & $\lambda A.b\equiv\lambda A'.b':\forall A.B$ \\
\texttt{forallEDF} & $A\equiv A':\mathrm{Sort}\,u$, $A{::}\Gamma\vdash B\equiv B':\mathrm{Sort}\,v$ & $\forall A.B\equiv\forall A'.B':\mathrm{Sort}(\mathrm{imax}\,uv)$ \\
\texttt{defeqDF} & $A\equiv B:\mathrm{Sort}\,u$, $e_1\equiv e_2:A$ & $e_1\equiv e_2:B$ \\
\texttt{beta} & $A{::}\Gamma\vdash e:B$, $a:A$ & $(\lambda A.e)\,a\equiv e[a]:B[a]$ \\
\texttt{eta} & $e:\forall A.B$ & $e\equiv\lambda A.(e{\uparrow})\,\mathrm{bvar}\,0:\forall A.B$ \\
\texttt{proofIrrel} & $p:\mathrm{Sort}\,0$, $h,h':p$ & $h\equiv h':p$ \\
\texttt{extra} & registered $\delta$-axiom, $ls$ level-wf & $e_1[ls]\equiv e_2[ls]:A[ls]$ \\
\texttt{pat} (iota) & Def.~\ref{rule:isdefeq-pat} & $e\equiv r_1.\mathrm{apply}\,m_1m_2:A$ \\
\bottomrule
\end{tabular}\par\smallskip
\end{center}
\end{definition}

\begin{definition}[The $\iota$ rule, \texttt{IsDefEq.pat}]
\label{rule:isdefeq-pat}
\contrib{iota}
\lean{Lean4Lean.VEnv.IsDefEq.pat}
\leanok
\uses{def:isdefeq, struct:venv, inductive:pattern-matches, def:pattern-check-realizes, inductive:pattern-rhs, def:venv-addpat}
\srcloc{Lean4Lean/Theory/Typing/Basic.lean}{60}
\thesisref{axioms.tex \S2.6.4, in the form of unique.tex ``Regularity of reductions'' item 1}

\lead{In short} A registered, matching, typed rule fires, giving \hl{the reduct at the
redex's own type}.
\begin{itemize}
\item \lead{Given} $(p,r)\in E.\texttt{pats}$; $e$ matches $p$ ($m_1,m_2$); $\Gamma\vdash e:A$;
  \texttt{chk} realizes $r_2$'s \texttt{Check}, each triple a definitional equality.
\item \lead{Then} $\Gamma \vdash e \equiv r_1.\mathrm{apply}\,m_1\,m_2 : A$.
\item \lead{Caveat} No premise types the reduct: subject reduction for $\iota$ is assumed,
  false over a merely \texttt{Ordered} environment (\ref{thm:wf-patsstrong};
  \srcloc{Lean4Lean/Theory/Typing/Basic.lean}{57}). The strong system's
  \texttt{IsDefEqStrong.pat} (\ref{inductive:isdefeqstrong-pat}) types the reduct instead.
\end{itemize}
\end{definition}

\begin{definition}[Derived judgments and well-formedness of environment entries]
\label{def:hastype-istype}
\lean{Lean4Lean.VEnv.HasType, Lean4Lean.VEnv.IsType, Lean4Lean.VEnv.IsDefEqU, Lean4Lean.VExpr.WF, Lean4Lean.VConstant.WF, Lean4Lean.VDefEq.WF}
\leanok
\uses{def:isdefeq}
\srcloc{Lean4Lean/Theory/Typing/Basic.lean}{68}
\thesisref{axioms.tex \S2.1 ($\Gamma \vdash \alpha\ \mathsf{type}$) and \S2.5}

\lead{In short} Five judgments on \texttt{IsDefEq}: typing, being a type, untyped equality,
\hl{well-formedness of a constant or an axiom}.
\begin{itemize}
\item \texttt{HasType}: the diagonal $e\equiv e:A$. \texttt{IsType}: $\exists u.\,A:\mathrm{Sort}\,u$.
\item \texttt{IsDefEqU}: \texttt{IsDefEq} with the type forgotten; \texttt{VExpr.WF} is its
  diagonal.
\item \texttt{VConstant.WF}: the constant's type is a type at its own arity, empty context.
\item \texttt{VDefEq.WF}: both axiom sides have its declared type there.
\end{itemize}
\end{definition}

\begin{definition}[Typing of a schematic reduction rule, \texttt{PatTyped}]
\label{def:pattyped}
\contrib{iota}
\lean{Lean4Lean.VEnv.PatTyped}
\leanok
\uses{inductive:pattern-matches, def:pattern-rhs-generic, def:hastype-istype, def:vlevel-inst}
\srcloc{Lean4Lean/Theory/Typing/Basic.lean}{92}
\thesisref{axioms.tex \S2.6.4, the parenthetical about substitution instances}

\lead{In short} A rule is typed when \hl{some generic instance} of redex and reduct share a
type.
\begin{itemize}
\item \lead{Given} a redex $e$ matching $p$ at \texttt{VLevel.params}\,$U$, $r_1$ generic in
  $m_2$ over $\Gamma$ (holes are exactly $\Gamma$'s variables), a type $B$.
\item \lead{Then} $\Gamma \vdash e : B$ and $\Gamma \vdash r_1.\mathrm{apply}(\ldots) : B$;
  the side conditions $r_2$ are not assumed.
\item \lead{Caveat} Existential in $U,\Gamma,e,m_2,B$; only the holes the reduct \emph{uses}
  are constrained (\srcloc{Lean4Lean/Theory/Typing/Pattern.lean}{181}): a rule can be typed
  at one convenient instantiation.
\end{itemize}
\end{definition}

\begin{definition}[Admissibility of a reduction rule, \texttt{PatWF}]
\label{def:patwf}
\contrib{iota}
\lean{Lean4Lean.VEnv.PatWF}
\leanok
\uses{def:pattyped, inductive:pattern-template-headed}
\srcloc{Lean4Lean/Theory/Typing/Basic.lean}{104}
\thesisref{axioms.tex \S2.6.4; no thesis counterpart as a separate predicate}

\lead{In short} Typed, and its reduct is a \hl{fixed closed $\lambda$-term applied to
arguments} (\texttt{TemplateHeaded}).
\begin{itemize}
\item Separates a \texttt{pats} entry from a definitional axiom: no instance of the reduct
  is a sort, a $\Pi$-type or a bare hole.
\item Monotone in the environment (\ref{fam:patwf-mono}); does not give subject reduction.
\end{itemize}
\end{definition}

\section{Well-formed environments}

An environment is built by a list of declaration steps, each checked before it is added.

\begin{definition}[Well-formedness of one declaration step, \texttt{VDecl.WF}]
\label{def:vdecl-wf}
\lean{Lean4Lean.VDecl.WF, Lean4Lean.VDefVal.WF, Lean4Lean.VEnv.addConsts, Lean4Lean.VEnv.addDefEqs}
\leanok
\uses{def:venv-basic, def:venv-addquot, def:ind-add-induct, def:hastype-istype, structure:ind-wf}
\srcloc{Lean4Lean/Theory/Typing/Env.lean}{18}
\thesisref{axioms.tex \S2.5, \S2.6, \S2.7.1}

\lead{In short} \texttt{VDecl.WF E d E'}: $d$ may be \hl{added to $E$}, in one of seven
forms.
\begin{center}
\begin{tabular}{lp{0.83\linewidth}}
\toprule
Form & Effect \\
\midrule
\texttt{axiom} & type is a type; added as a constant. \\
\texttt{def} & value typed at the declared type; add constant, then its $\delta$ rule. \\
\texttt{opaque} & same premise; only the constant is added, value hidden. \\
\texttt{example} & same premise; environment unchanged. \\
\texttt{mutualDef} & block types checked, added (\texttt{addConsts}); values typed in the
  extension; every $\delta$ rule added (\texttt{addDefEqs}). \\
\texttt{quot} & Definition~\ref{def:venv-addquot}. \\
\texttt{induct} & Definition~\ref{structure:ind-wf}. \\
\bottomrule
\end{tabular}\par\smallskip
\end{center}
\end{definition}

\begin{definition}[Well-formed environment, \texttt{VEnv.WF}]
\label{def:venv-wf}
\lean{Lean4Lean.VEnv.WF', Lean4Lean.VEnv.WF}
\leanok
\uses{def:vdecl-wf, def:venv-basic}
\srcloc{Lean4Lean/Theory/Typing/Env.lean}{48}

\lead{In short} \texttt{WF' ds E}: $E$ is reached from empty by \texttt{ds}, each step
well-formed; \texttt{VEnv.WF E} is its \hl{existential closure}.
\begin{itemize}
\item The formal notion of an environment the kernel could have built.
\end{itemize}
\end{definition}

\begin{definition}[Well-formed prefix, \texttt{WFPrefix}]
\label{def:wfprefix}
\contrib{iota}
\lean{Lean4Lean.VEnv.WFPrefix}
\leanok
\uses{def:vdecl-wf}
\srcloc{Lean4Lean/Theory/Typing/Env.lean}{57}
\thesisref{no thesis counterpart; bookkeeping for the formal induction}

\lead{In short} \texttt{E.WFPrefix}\,$E_0$: $E$ extends $E_0$ by a run of \hl{well-formed
steps} (constructors \texttt{rfl}, \texttt{decl}).
\begin{itemize}
\item What the strengthening argument of \ref{thm:wf-strong} walks through.
\end{itemize}
\end{definition}

\begin{lemma}[Prefix of a prefix, \texttt{WFPrefix.trans}]
\label{thm:wfprefix-trans}
\contrib{iota}
\lean{Lean4Lean.VEnv.WFPrefix.trans}
\leanok
\uses{def:wfprefix}
\srcloc{Lean4Lean/Theory/Typing/Env.lean}{62}
\thesisref{no thesis counterpart}

\lead{In short} \texttt{WFPrefix} is \hl{transitive}: \texttt{E.WFPrefix}\,$E_1$ and $E_1$
\texttt{.WFPrefix}\,$E_2$ give \texttt{E.WFPrefix}\,$E_2$.
\end{lemma}
\begin{proof}
\uses{def:wfprefix}
\leanok
\lead{Idea} Induction: \texttt{rfl} returns the second hypothesis, \texttt{decl} rebuilds
around it.
\end{proof}

\section{Contexts and the ordered invariant}

\texttt{Ordered} is the builds-in-order invariant the structural lemmas need.

\begin{lemma}[Context liftings and instantiations]
\label{fam:ctx-lift}
\lean{Lean4Lean.Ctx.LiftN, Lean4Lean.Ctx.Lift', Lean4Lean.Ctx.InstN, Lean4Lean.Ctx.LiftN.one, Lean4Lean.Ctx.LiftN.isSuffix, Lean4Lean.Ctx.LiftN.comp, Lean4Lean.Ctx.LiftN.right, Lean4Lean.Ctx.liftN_iff_lift', Lean4Lean.Ctx.Lift'.comp, Lean4Lean.Ctx.Lift'.of_cons_skip, Lean4Lean.Ctx.Lift'.depth_zero}
\leanok
\uses{def:liftn, def:lift-prime, def:inst}
\srcloc{Lean4Lean/Theory/Typing/Lemmas.lean}{8}
\thesisref{typesys.tex thm:weak and thm:subst, the context side}

\lead{In short} Context-level \hl{lifting and instantiation}: \texttt{LiftN} inserts $n$
entries $k$ deep, \texttt{Lift'} along a general lift word, \texttt{InstN} substitutes at
depth $k$.
\begin{itemize}
\item Family: one-entry case, suffix characterisation, composition, the two presentations'
  equivalence, weakening a closed context on the right.
\end{itemize}
\end{lemma}
\begin{proof}
\leanok
\lead{Idea} Routine inductions on lifting and instantiation.
\end{proof}

\begin{lemma}[Structural lemmas for variable lookup]
\label{fam:lookup-struct}
\lean{Lean4Lean.Lookup.lt, Lean4Lean.Lookup.ofLt, Lean4Lean.Lookup.uniq, Lean4Lean.Lookup.weak', Lean4Lean.Lookup.weakN, Lean4Lean.Lookup.weakU_inv, Lean4Lean.Lookup.weak'_iff, Lean4Lean.Lookup.weakN_iff, Lean4Lean.Lookup.instL}
\leanok
\uses{def:lookup, fam:ctx-lift}
\srcloc{Lean4Lean/Theory/Typing/Lemmas.lean}{90}
\thesisref{typesys.tex thm:weak}

\lead{In short} \texttt{Lookup} is \hl{well-behaved}: bounded, total below the context
length, functional, preserved/reflected by lifting, commutes with level instantiation.
\end{lemma}
\begin{proof}
\uses{fam:ctx-lift}
\leanok
\lead{Idea} Inductions on \texttt{Lookup} and on lifting.
\end{proof}

\begin{lemma}[Pointwise predicates on contexts]
\label{def:onctx}
\lean{Lean4Lean.OnCtx, Lean4Lean.OnCtx.lookup, Lean4Lean.OnCtx.mono, Lean4Lean.CtxClosed, Lean4Lean.CtxClosed.lookup}
\leanok
\uses{def:lookup}
\srcloc{Lean4Lean/Theory/Typing/Lemmas.lean}{152}
\thesisref{axioms.tex \S2.1, the judgment $\vdash \Gamma\ \mathsf{ok}$}

\lead{In short} \texttt{OnCtx}\,$\Gamma P$: $P$ of \hl{every entry} and the context to its
right; \texttt{OnCtx}\,$\Gamma$\,(\texttt{IsType}\,$U$) is $\vdash \Gamma\ \mathsf{ok}$.
\begin{itemize}
\item \texttt{CtxClosed}: specialised to closedness; transfers $P$ along \texttt{Lookup}.
\end{itemize}
\end{lemma}
\begin{proof}
\leanok
\lead{Idea} Structural recursion; \texttt{OnCtx} is recursive, so each case is a two-way
\texttt{match}.
\end{proof}

\begin{lemma}[Registering a rule: basic facts]
\label{fam:addpat}
\contrib{iota}
\lean{Lean4Lean.VEnv.addPat_le, Lean4Lean.VEnv.addPat_self}
\leanok
\uses{def:venv-addpat, struct:venv-le}
\srcloc{Lean4Lean/Theory/Typing/Lemmas.lean}{198}
\thesisref{no thesis counterpart}

\lead{In short} Registering a rule only \hl{grows the environment} ($E \le E.\texttt{addPat}
\,p\,r$), and the rule is present afterwards.
\end{lemma}
\begin{proof}
\uses{def:venv-addpat, struct:venv-le}
\leanok
\lead{Idea} One-liners from the definitions; the change also gave \texttt{addConst\_le} and
\texttt{addDefEq\_le} a third component (\texttt{VEnv.LE} gained a \texttt{pats} field).
\end{proof}

\begin{definition}[Ordered environment, \texttt{Ordered}]
\label{def:ordered}
\contrib{mixed}
\lean{Lean4Lean.VEnv.Ordered}
\leanok
\uses{def:venv-basic, def:venv-addpat, def:hastype-istype, def:patwf}
\srcloc{Lean4Lean/Theory/Typing/Lemmas.lean}{257}
\thesisref{no thesis counterpart; the formal stand-in for ``built in order''}

\lead{In short} Built from empty by steps each \hl{checked in the environment before it}.
\begin{itemize}
\item \texttt{const}: type is a type. \texttt{defeq}: both axiom sides typed. \texttt{pat}
  (iota, \srcloc{Lean4Lean/Theory/Typing/Lemmas.lean}{265}): satisfies \texttt{PatWF}.
\item \lead{Caveat} Alone, does not make \texttt{IsDefEq} well-behaved: \texttt{defeq} admits
  any well-typed axiom, letting a registered $\iota$ rule take a well-typed redex to an
  ill-typed reduct, while \ref{rule:isdefeq-pat} forces the reduct typed regardless.
\end{itemize}
\end{definition}

\begin{definition}[A predicate on every constant type and axiom side, \texttt{OnTypes}]
\label{def:ontypes}
\lean{Lean4Lean.VEnv.OnTypes, Lean4Lean.VEnv.OnTypes.mono}
\leanok
\uses{struct:venv}
\srcloc{Lean4Lean/Theory/Typing/Lemmas.lean}{267}

\lead{In short} Every constant type and axiom side of $E$ satisfies $P$: the shape of every
\hl{``environment is good''} statement here.
\begin{itemize}
\item Antitone in $E$, monotone in $P$.
\item \lead{Caveat} Says nothing about \texttt{pats}: why the $\iota$ obligation could not
  ride \ref{thm:ordered-induction}, forcing \ref{thm:wf-strong}'s bootstrap from scratch.
\end{itemize}
\end{definition}

\begin{theorem}[Induction over an ordered environment]
\label{thm:ordered-induction}
\contrib{mixed}
\lean{Lean4Lean.VEnv.Ordered.induction}
\leanok
\uses{def:ordered, def:ontypes}
\srcloc{Lean4Lean/Theory/Typing/Lemmas.lean}{275}

\lead{In short} Proves \hl{\texttt{OnTypes E (motive E)}} of an ordered $E$.
\begin{itemize}
\item \lead{Given} the motive is monotone and holds of every closed $[\,] \vdash e : A$ typed
  where the environment's own entries already satisfy it.
\item \lead{Then} \texttt{OnTypes E (motive E)}.
\end{itemize}
\end{theorem}
\begin{proof}
\uses{def:ordered, def:ontypes, fam:addpat}
\leanok
\lead{Idea} Induction on \texttt{Ordered}: \texttt{const}/\texttt{defeq} unfold and split;
the \texttt{pat} case (iota) transports along \texttt{addPat\_le}, contributing nothing.
\end{proof}

\section{Structural metatheory}

The core lemmas, mostly over an ordered environment: regularity, monotonicity, weakening,
substitution and inversion.

\begin{lemma}[Regularity: free variables are in the context]
\label{thm:isdefeq-closedn}
\contrib{mixed}
\lean{Lean4Lean.VEnv.IsDefEq.closedN', Lean4Lean.VEnv.Ordered.closed, Lean4Lean.VEnv.Ordered.closedC, Lean4Lean.VEnv.IsDefEq.closedN, Lean4Lean.VExpr.WF.closedN}
\leanok
\uses{def:isdefeq, def:onctx, def:ontypes}
\srcloc{Lean4Lean/Theory/Typing/Lemmas.lean}{329}
\thesisref{typesys.tex Lemma (Regularity) thm:reg item 2}

\lead{In short} A closed environment and context keep a derivation's terms \hl{closed below
$|\Gamma|$}.
\begin{itemize}
\item \texttt{Ordered.closed} discharges the environment side; \texttt{closedC} specialises
  to constant types.
\end{itemize}
\end{lemma}
\begin{proof}
\uses{def:isdefeq, thm:ordered-induction, family:pattern-transport-iota}
\leanok
\lead{Idea} Induction using closedness of substitution/lifting; the \texttt{pat} case (iota)
gets the reduct's closedness from \texttt{Pattern.RHS.apply\_closedN}.
\end{proof}

\begin{theorem}[Monotonicity in the environment]
\label{thm:isdefeq-mono}
\contrib{mixed}
\lean{Lean4Lean.VEnv.IsDefEq.mono, Lean4Lean.VEnv.HasType.mono, Lean4Lean.VEnv.IsType.mono, Lean4Lean.VEnv.IsDefEqU.mono, Lean4Lean.VExpr.WF.mono, Lean4Lean.VConstant.WF.mono, Lean4Lean.VDefEq.WF.mono}
\leanok
\uses{def:isdefeq, struct:venv-le}
\srcloc{Lean4Lean/Theory/Typing/Lemmas.lean}{397}

\lead{In short} Every judgment derivable in $E$ is derivable in any \hl{$E' \ge E$}, and
likewise for the derived judgments.
\end{theorem}
\begin{proof}
\uses{def:isdefeq, struct:venv-le}
\leanok
\lead{Idea} Induction replaying each rule; \texttt{constDF}, \texttt{extra}, \texttt{pat}
(iota) use the three components of \texttt{VEnv.LE}.
\end{proof}

\begin{lemma}[Entries of an ordered environment are well-formed in it]
\label{fam:ordered-entrywf}
\contrib{mixed}
\lean{Lean4Lean.VEnv.Ordered.constWF, Lean4Lean.VEnv.Ordered.defEqWF}
\leanok
\uses{def:ordered, def:hastype-istype}
\srcloc{Lean4Lean/Theory/Typing/Lemmas.lean}{439}

\lead{In short} Every constant of an ordered environment has a \hl{well-formed type} there;
every axiom has both sides typed there.
\end{lemma}
\begin{proof}
\uses{def:ordered, thm:isdefeq-mono, fam:addpat}
\leanok
\lead{Idea} Induction on \texttt{Ordered}, transporting along each step via
\texttt{VConstant.WF.mono}/\texttt{VDefEq.WF.mono}; \texttt{pat} (iota) via
\texttt{addPat\_le}.
\end{proof}

\begin{lemma}[Rule well-formedness is monotone]
\label{fam:patwf-mono}
\contrib{iota}
\lean{Lean4Lean.VEnv.PatTyped.mono, Lean4Lean.VEnv.PatWF.mono}
\leanok
\uses{def:pattyped, def:patwf, struct:venv-le}
\srcloc{Lean4Lean/Theory/Typing/Lemmas.lean}{464}
\thesisref{no thesis counterpart}

\lead{In short} \texttt{PatTyped} and \texttt{PatWF} transfer along \hl{$E \le E'$}.
\end{lemma}
\begin{proof}
\uses{def:pattyped, def:patwf, thm:isdefeq-mono}
\leanok
\lead{Idea} Move the two typing judgments by \texttt{HasType.mono}; \texttt{TemplateHeaded}
does not mention the environment. What makes \texttt{Ordered.pat} legitimate.
\end{proof}

\begin{theorem}[Registered rules of an ordered environment are well-formed]
\label{thm:ordered-patwf}
\contrib{iota}
\lean{Lean4Lean.VEnv.Ordered.patWF, Lean4Lean.VEnv.Ordered.patHeaded}
\leanok
\uses{def:ordered, def:patwf, inductive:pattern-template-headed}
\srcloc{Lean4Lean/Theory/Typing/Lemmas.lean}{474}
\thesisref{no thesis counterpart; the \texttt{pats} analogue of \ref{fam:ordered-entrywf}}

\lead{In short} Every rule in an ordered environment's \texttt{pats} satisfies
\texttt{PatWF}; its reduct is \hl{template-headed} (\texttt{patHeaded}).
\end{theorem}
\begin{proof}
\uses{def:ordered, fam:patwf-mono, fam:addpat}
\leanok
\lead{Idea} Induction on \texttt{Ordered}: the rule was just added (hypothesis supplies
\texttt{PatWF}) or transports via \texttt{PatWF.mono}.
\end{proof}

\begin{lemma}[Levels occurring in a derivation are well-formed]
\label{thm:isdefeq-levelwf}
\contrib{mixed}
\lean{Lean4Lean.VEnv.IsDefEq.levelWF}
\leanok
\uses{def:isdefeq, def:vlevel-wf, def:onctx}
\srcloc{Lean4Lean/Theory/Typing/Lemmas.lean}{498}

\lead{In short} If $\Gamma$'s types mention only levels below $U$, a derivation of
$\Gamma \vdash e_1 \equiv e_2 : A$ shows the same of $e_1,e_2,A$ -- \hl{no environment
hypothesis needed}.
\end{lemma}
\begin{proof}
\uses{def:isdefeq, family:pattern-transport-iota}
\leanok
\lead{Idea} Induction; \texttt{pat} (iota) via \texttt{Pattern.RHS.apply\_levelWF} fed by
\texttt{Matches.levelWF}.
\end{proof}

\begin{theorem}[Weakening]
\label{thm:isdefeq-weakn}
\contrib{mixed}
\lean{Lean4Lean.VEnv.IsDefEq.weakN, Lean4Lean.VEnv.HasType.weakN, Lean4Lean.VEnv.IsType.weakN, Lean4Lean.VEnv.IsDefEq.weak, Lean4Lean.VEnv.IsDefEq.weakR, Lean4Lean.VEnv.IsDefEq.weak0, Lean4Lean.VEnv.IsDefEq.weak'}
\leanok
\uses{def:isdefeq, fam:ctx-lift, def:ordered}
\srcloc{Lean4Lean/Theory/Typing/Lemmas.lean}{546}
\thesisref{typesys.tex Lemma (Weakening) thm:weak items 1--2}

\lead{In short} Over an ordered environment, a derivation transports along any \hl{context
insertion}: $\Gamma' \vdash e_1{\uparrow} \equiv e_2{\uparrow} : A{\uparrow}$ when $\Gamma'$
inserts $n$ entries $k$ deep into $\Gamma$.
\begin{itemize}
\item Corollaries: one binder, a closed suffix on the right, from empty, a general lift word.
\end{itemize}
\end{theorem}
\begin{proof}
\uses{def:isdefeq, thm:isdefeq-closedn, family:pattern-transport-master, family:pattern-transport-iota}
\leanok
\lead{Idea} Induction commuting lifting past instantiation; \texttt{pat} (iota) rewrites by
\texttt{Pattern.RHS.liftN\_apply} and re-derives the match.
\end{proof}

\begin{lemma}[Universe-level instantiation]
\label{thm:isdefeq-instl}
\contrib{mixed}
\lean{Lean4Lean.VEnv.IsDefEq.instL, Lean4Lean.VEnv.HasType.instL, Lean4Lean.VEnv.IsType.instL, Lean4Lean.VEnv.IsDefEqU.instL, Lean4Lean.OnCtx.instL}
\leanok
\uses{def:isdefeq, def:vlevel-inst}
\srcloc{Lean4Lean/Theory/Typing/Lemmas.lean}{646}
\thesisref{axioms.tex \S2.1, universe-polymorphic constants}

\lead{In short} A derivation at arity $U$ instantiates to any level list well-formed at
$U'$, with \hl{context, terms and type instantiated}.
\end{lemma}
\begin{proof}
\uses{def:isdefeq, family:pattern-transport-iota}
\leanok
\lead{Idea} Induction, levels substituted pointwise; \texttt{pat} (iota) via
\texttt{Pattern.RHS.instL\_apply}.
\end{proof}

\begin{theorem}[Substitution]
\label{thm:isdefeq-instn}
\contrib{mixed}
\lean{Lean4Lean.VEnv.IsDefEq.instN, Lean4Lean.VEnv.HasType.instN, Lean4Lean.VEnv.IsType.instN, Lean4Lean.VEnv.IsDefEqU.instN, Lean4Lean.Ctx.InstN.wf}
\leanok
\uses{def:isdefeq, fam:ctx-lift, def:ordered}
\srcloc{Lean4Lean/Theory/Typing/Lemmas.lean}{691}
\thesisref{typesys.tex Lemma (Properties of substitution) thm:subst item subst\_ty}

\lead{In short} Over an ordered environment, \hl{substituting a well-typed term} transports
a derivation: $\Gamma_0 \vdash e_0 : A_0$ and $\Gamma_1 \vdash e_1 \equiv e_2 : A$ give
$\Gamma \vdash e_1[e_0] \equiv e_2[e_0] : A[e_0]$ (\texttt{InstN} at depth $k$).
\end{theorem}
\begin{proof}
\uses{def:isdefeq, thm:isdefeq-weakn, family:pattern-transport-master, family:pattern-transport-iota}
\leanok
\lead{Idea} Induction with an inner induction on \texttt{Ctx.InstN}; \texttt{pat} (iota) via
\texttt{Pattern.RHS.instN\_apply} -- where the rule's substitution-closure is discharged.
\end{proof}

\begin{lemma}[Context conversion]
\label{fam:defeqdfc}
\lean{Lean4Lean.VEnv.IsDefEqCtx, Lean4Lean.VEnv.IsDefEq.defeqDF_l', Lean4Lean.VEnv.IsDefEq.defeqDF_l, Lean4Lean.VEnv.IsDefEq.defeqDFC', Lean4Lean.VEnv.IsDefEq.defeqDFC, Lean4Lean.VEnv.HasType.defeqDFC, Lean4Lean.VEnv.IsType.defeqDFC, Lean4Lean.VEnv.IsDefEqU.defeqDFC, Lean4Lean.VEnv.IsDefEqCtx.symm, Lean4Lean.VEnv.IsDefEqCtx.isType, Lean4Lean.VEnv.IsDefEqCtx.closed, Lean4Lean.VEnv.IsDefEqCtx.refl, Lean4Lean.VEnv.IsDefEqCtx.length_eq, Lean4Lean.VEnv.IsDefEqCtx.isSuffix}
\leanok
\uses{def:isdefeq, def:onctx}
\srcloc{Lean4Lean/Theory/Typing/Lemmas.lean}{302}
\thesisref{not stated separately; an admissible rule the de Bruijn presentation needs}

\lead{In short} \texttt{IsDefEqCtx} relates two contexts sharing a suffix, \hl{pointwise
definitionally equal} above it; derivations transport along it.
\end{lemma}
\begin{proof}
\uses{thm:isdefeq-weakn, thm:isdefeq-instn, fam:ctx-lift}
\leanok
\lead{Idea} The one-binder case weakens then substitutes the zero variable back; the general
case iterates. Untouched by either branch.
\end{proof}

\begin{theorem}[Inversion for $\Pi$-types]
\label{thm:isdefeq-foralle-inv}
\contrib{mixed}
\lean{Lean4Lean.VEnv.IsDefEq.forallE_inv', Lean4Lean.VEnv.HasType.forallE_inv, Lean4Lean.VEnv.IsType.forallE_inv}
\leanok
\uses{def:isdefeq, def:ordered, def:hastype-istype}
\srcloc{Lean4Lean/Theory/Typing/Lemmas.lean}{804}
\thesisref{typesys.tex thm:reg; a weak form of unique.tex thm:1dinv item 2}

\lead{In short} Over an ordered environment, a $\Pi$-type on one side forces \hl{both domain
and codomain to be types}.
\end{theorem}
\begin{proof}
\uses{thm:ordered-induction, thm:isdefeq-instn, thm:isdefeq-weakn, thm:ordered-patwf, inductive:pattern-template-headed}
\leanok
\lead{Idea} Induction with the side as a disjunction; \texttt{pat} (iota) is discharged
outright by \texttt{TemplateHeaded.apply\_ne\_forallE} -- else this would need the open
\ref{thm:injectivity-open}.
\end{proof}

\begin{lemma}[Inversion for sorts]
\label{thm:isdefeq-sort-inv}
\contrib{mixed}
\lean{Lean4Lean.VEnv.IsDefEq.sort_inv', Lean4Lean.VEnv.IsDefEq.sort_inv_l, Lean4Lean.VEnv.IsDefEq.sort_inv_r, Lean4Lean.VEnv.HasType.sort_inv, Lean4Lean.VEnv.IsType.sort_inv, Lean4Lean.VEnv.IsDefEq.sort_r}
\leanok
\uses{def:isdefeq, def:ordered, def:vlevel-wf}
\srcloc{Lean4Lean/Theory/Typing/Lemmas.lean}{866}
\thesisref{unique.tex thm:1dinv item 1, weak form}

\lead{In short} Over an ordered environment, $\mathrm{Sort}\,u$ on one side makes $u$ a
\hl{well-formed level}; does not compare two sorts.
\end{lemma}
\begin{proof}
\uses{fam:ordered-entrywf, thm:ordered-patwf, inductive:pattern-template-headed}
\leanok
\lead{Idea} Mirrors the $\Pi$ case; \texttt{pat} (iota) via
\texttt{TemplateHeaded.apply\_ne\_sort}.
\end{proof}

\begin{theorem}[Regularity: the type of a judgment is a type]
\label{thm:isdefeq-istype}
\contrib{mixed}
\lean{Lean4Lean.VEnv.IsDefEq.isType', Lean4Lean.VEnv.Ordered.isType, Lean4Lean.VEnv.IsDefEq.isType}
\leanok
\uses{def:isdefeq, def:ordered, def:onctx, def:hastype-istype}
\srcloc{Lean4Lean/Theory/Typing/Lemmas.lean}{912}
\thesisref{typesys.tex Lemma (Regularity) thm:reg item 4}

\lead{In short} In an ordered environment with a well-formed context, a derivation's type is
itself \hl{a type}.
\end{theorem}
\begin{proof}
\uses{thm:isdefeq-foralle-inv, thm:isdefeq-sort-inv, thm:isdefeq-instn, thm:ordered-induction}
\leanok
\lead{Idea} Application inverts the $\Pi$ and substitutes; constant/axiom read off
\texttt{Ordered.isType}; \texttt{pat} (iota) reuses the induction hypothesis for $e:A$.
\end{proof}

\begin{lemma}[Substitution of definitionally equal arguments]
\label{thm:isdefeq-instdf}
\lean{Lean4Lean.VEnv.IsDefEq.instDF}
\leanok
\uses{def:isdefeq, def:onctx}
\srcloc{Lean4Lean/Theory/Typing/Lemmas.lean}{952}
\thesisref{typesys.tex thm:subst item 3}

\lead{In short} Over an ordered environment and well-formed context, $A :: \Gamma \vdash f
\equiv f' : B$ and $\Gamma \vdash a \equiv a' : A$ give \hl{$\Gamma \vdash f[a] \equiv
f'[a'] : B[a]$}.
\end{lemma}
\begin{proof}
\uses{thm:isdefeq-istype, thm:isdefeq-sort-inv}
\leanok
\lead{Idea} Rewrite $f[a]$ as $(\lambda x{:}A.f)\,a$ by $\beta$, apply congruence, convert
back; used for the term and, via \ref{thm:isdefeq-istype}, the type.
\end{proof}

\begin{lemma}[Pointwise definitionally equal substitutions]
\label{fam:substeq}
\lean{Lean4Lean.VEnv.Ctx.SubstEq, Lean4Lean.VEnv.Ctx.SubstEq.left, Lean4Lean.VEnv.Ctx.SubstEq.lookup, Lean4Lean.VEnv.Ctx.SubstEq.skip, Lean4Lean.VEnv.Ctx.SubstEq.lift, Lean4Lean.VEnv.Ctx.SubstEq.lift_at, Lean4Lean.VEnv.Ctx.SubstEq.wf, Lean4Lean.VEnv.Ctx.SubstEq.id}
\leanok
\uses{def:isdefeq, def:subst, def:onctx}
\srcloc{Lean4Lean/Theory/Typing/Lemmas.lean}{973}
\thesisref{the simultaneous-substitution generalisation of typesys.tex thm:subst}

\lead{In short} \texttt{Ctx.SubstEq}: a pair of substitutions to $\Gamma_0$, \hl{pointwise
definitionally equal} at the corresponding types.
\begin{itemize}
\item Diagonal projection, lookup, codomain weakening, extension under a binder
  (\texttt{lift\_at}), source well-formedness, identity. Consumed by
  \texttt{IsDefEq.substDF} (Chapter~\ref{chap:metatheory}).
\end{itemize}
\end{lemma}
\begin{proof}
\uses{thm:isdefeq-weakn, def:subst}
\leanok
\lead{Idea} Inductions with \texttt{VExpr.Subst} bookkeeping, via \texttt{IsDefEq.weak'}.
All upstream.
\end{proof}

\section{Tactic support and the quotient block}

Small tactic support closes typing goals for the quotient block: four constants and one
$\delta$ rule.

\begin{lemma}[Tactic support for closed typing goals]
\label{fam:typing-meta}
\lean{Lean4Lean.VEnv.with_addConst, Lean4Lean.VEnv.Lookup.zero', Lean4Lean.VEnv.Lookup.succ', Lean4Lean.VEnv.HasType.app', Lean4Lean.VEnv.HasType.const'}
\leanok
\uses{def:lookup, def:hastype-istype, def:venv-basic, def:ordered}
\srcloc{Lean4Lean/Theory/Typing/Meta.lean}{6}

\lead{In short} Up-to-equation lookup/typing variants chained into \hl{tactics for closed
typing goals} (\texttt{lookup\_tac}, \texttt{type\_tac}), plus \texttt{with\_addConst}.
\begin{itemize}
\item Used only on the quotient block: twelve sites in \texttt{QuotLemmas.lean}, a hundred
  in the contributed \srcloc{Lean4Lean/Verify/Environment/Quot.lean}{1}.
\end{itemize}
\end{lemma}
\begin{proof}
\leanok
\lead{Idea} Variants are \texttt{eq $\triangleright$ rule}; the macros chain them;
\texttt{with\_addConst} peels one \texttt{Option.bind}.
\end{proof}

\begin{lemma}[Adding the quotient constants preserves orderedness]
\label{thm:addquot-wf}
\lean{Lean4Lean.VEnv.addQuot_WF}
\leanok
\uses{def:ordered, def:venv-addquot}
\srcloc{Lean4Lean/Theory/Typing/QuotLemmas.lean}{7}
\thesisref{axioms.tex \S2.7.1 (quotient types)}

\lead{In short} Adding \texttt{Quot}, \texttt{Quot.mk/lift/ind} and the computation rule
keeps the environment \hl{ordered}, given $E$ ordered and \texttt{QuotReady}.
\begin{itemize}
\item \lead{Caveat} The computation rule is a \texttt{VDefEq} ($\delta$ rule), not a
  \texttt{pats} entry, though the thesis calls it $\iota$; \ref{thm:wf-patsstrong} must
  handle this asymmetry.
\end{itemize}
\end{lemma}
\begin{proof}
\uses{fam:typing-meta, def:ordered, def:venv-addquot}
\leanok
\lead{Idea} Four nested \texttt{with\_addConst}, each goal closed by \texttt{type\_tac},
ending with \texttt{Ordered.defeq}.
\end{proof}

\begin{lemma}[\texttt{addQuot} as an explicit chain of five steps]
\label{thm:addquot-chain}
\contrib{iota}
\lean{Lean4Lean.VEnv.addQuot_chain}
\leanok
\uses{def:venv-addquot, def:venv-basic, def:hastype-istype}
\srcloc{Lean4Lean/Theory/Typing/QuotLemmas.lean}{19}
\thesisref{axioms.tex \S2.7.1}

\lead{In short} Unpacks \texttt{addQuot} into its \hl{four intermediate environments}, each
constant's well-formedness, and \texttt{quotDefEq.WF} at the last.
\end{lemma}
\begin{proof}
\uses{def:venv-addquot, fam:typing-meta}
\leanok
\lead{Idea} Split the four binds, build the \texttt{HasObjects} chain, close each goal by
\texttt{type\_tac}; feeds \ref{thm:addquot-strong}.
\end{proof}

\begin{lemma}[Recovering the quotient objects from the extended environment]
\label{fam:addquot-objs}
\lean{Lean4Lean.VEnv.addQuot_objs, Lean4Lean.VEnv.addQuot_quot, Lean4Lean.VEnv.addQuot_quotMk, Lean4Lean.VEnv.addQuot_quotLift, Lean4Lean.VEnv.addQuot_quotInd, Lean4Lean.VEnv.addQuot_defeq}
\leanok
\uses{def:venv-addquot}
\srcloc{Lean4Lean/Theory/Typing/QuotLemmas.lean}{46}
\thesisref{axioms.tex \S2.7.1}

\lead{In short} After \texttt{addQuot}, the environment holds \hl{the four quotient
constants} at their declared types and the quotient axiom.
\end{lemma}
\begin{proof}
\leanok
\lead{Idea} Peel the four binds, then project.
\end{proof}

\section{From well-formedness to the ordered invariant}

\texttt{VEnv.WF} is what the kernel guarantees; \texttt{Ordered} is what the lemmas need.

\begin{lemma}[Bookkeeping for blocks of constants and definitional equations]
\label{fam:addconsts}
\lean{Lean4Lean.VEnv.addConsts_le, Lean4Lean.VEnv.addConst_eq_none, Lean4Lean.VEnv.addConst_constants_eq, Lean4Lean.VEnv.exists_addConsts, Lean4Lean.VEnv.addConsts_congr, Lean4Lean.VEnv.addConsts_ordered, Lean4Lean.VEnv.addConsts_constants, Lean4Lean.VEnv.addDefEqs_ordered}
\leanok
\uses{def:vdecl-wf, def:ordered, def:venv-basic}
\srcloc{Lean4Lean/Theory/Typing/EnvLemmas.lean}{8}
\thesisref{axioms.tex \S2.5 (mutual definitions)}

\lead{In short} Adding a constant block \hl{grows the environment}, succeeds when names are
fresh and distinct, and preserves orderedness when every constant is typed already.
\begin{itemize}
\item Only fresh/distinct $\Rightarrow$ succeeds is proved; adding the block's $\delta$
  rules afterwards (\texttt{addDefEqs}) also preserves orderedness.
\end{itemize}
\end{lemma}
\begin{proof}
\uses{def:ordered, def:vdecl-wf}
\leanok
\lead{Idea} Inductions unfolding \texttt{foldlM} one step at a time.
\end{proof}

\begin{theorem}[A well-formed environment is ordered]
\label{thm:wf-ordered}
\lean{Lean4Lean.VEnv.WF.ordered, Lean4Lean.instCoeOutWFOrdered}
\leanok
\uses{def:venv-wf, def:ordered}
\srcloc{Lean4Lean/Theory/Typing/EnvLemmas.lean}{87}

\lead{In short} Every well-formed environment satisfies \hl{\texttt{Ordered}}; a
\texttt{CoeOut} instance makes the passage automatic.
\end{theorem}
\begin{proof}
\uses{def:venv-wf, fam:addconsts, thm:addquot-wf, thm:indlem-add-induct-wf}
\leanok
\lead{Idea} Induction on \texttt{VEnv.WF'}, one case per form; on upstream \texttt{induct} was
vacuous (three \texttt{sorry}s), the iota branch gives it real content.
\end{proof}

\section{The strong-system bootstrap}

Chapter~\ref{chap:metatheory} needs every constant and axiom of a well-formed environment
strongly typed. Upstream's twelve-line entry point does not survive $\iota$, since
\ref{def:ontypes} ignores \texttt{pats}; this section rebuilds it.

\begin{lemma}[Declaration steps and prefixes only grow the environment]
\label{thm:vdecl-wf-le}
\contrib{iota}
\lean{Lean4Lean.VDecl.WF.le, Lean4Lean.VEnv.WFPrefix.le}
\leanok
\uses{def:vdecl-wf, def:wfprefix, struct:venv-le}
\srcloc{Lean4Lean/Theory/Typing/EnvLemmas.lean}{109}
\thesisref{no thesis counterpart}

\lead{In short} Every \texttt{VDecl.WF} step satisfies \hl{$E \le E'$}; hence
\texttt{E.WFPrefix}\,$E_0$ implies $E_0 \le E$.
\end{lemma}
\begin{proof}
\uses{fam:addconsts, family:indlem-pats-defeqs-preserved, family:indlem-stage-le}
\leanok
\lead{Idea} Case analysis composing the per-operation \texttt{le} lemmas.
\end{proof}

\begin{definition}[Subject reduction of $\iota$ along the strengthening argument, \texttt{PatsStrong}]
\label{def:patsstrong}
\contrib{iota}
\lean{Lean4Lean.VEnv.PatsStrong}
\leanok
\uses{def:wfprefix, def:patsstrongon, def:ordered, struct:venv-le}
\srcloc{Lean4Lean/Theory/Typing/EnvLemmas.lean}{130}
\thesisref{unique.tex ``Regularity of reductions'' item 1; typesys.tex ``Regularity continued'' item red\_equiv}

\lead{In short} \texttt{E.PatsStrong}: every ordered $E_0$ reached from a well-formed prefix
of $E$ by a \hl{constant-only extension} has subject reduction for its rules.
\begin{itemize}
\item \lead{Given} \texttt{E.WFPrefix}\,$E_1$, $E_1 \le E_0 \le E$, matching
  \texttt{defeqs}/\texttt{pats}, $E_0$ ordered.
\item \lead{Then} \texttt{PatsStrongOn}\,$E_0$
  (\srcloc{Lean4Lean/Theory/Typing/EnvLemmas.lean}{279}) -- exactly the stage environments of
  an inductive block and the intermediate environments of \texttt{addQuot}.
\item \lead{Caveat} Discharged only by \texttt{sorry} (\ref{thm:wf-patsstrong}): unverified
  whether provable at this generality, since $E_1$ may already carry $\delta$ rules and the
  quotient rule.
\end{itemize}
\end{definition}

\begin{theorem}[Strong typing survives a fold of constant additions]
\label{thm:foldlm-addconst-strong}
\contrib{iota}
\lean{Lean4Lean.VEnv.foldlM_addConst_strong}
\leanok
\uses{def:envstrong, def:patsstrongon, def:ontypes, def:ordered}
\srcloc{Lean4Lean/Theory/Typing/EnvLemmas.lean}{137}
\thesisref{no thesis counterpart; machinery for the annotated judgment}

\lead{In short} Folding \texttt{addConst} \hl{preserves strong typing} when every added
constant is typed as it is added and every intermediate stage has subject reduction.
\end{theorem}
\begin{proof}
\uses{thm:ontypes-addconst, family:indlem-pats-defeqs-preserved, family:indlem-foldlm}
\leanok
\lead{Idea} Induction applying \texttt{OnTypes.addConst}, re-aiming \texttt{PatsStrongOn} at
each stage; serves \texttt{addConsts} and \texttt{addInduct}'s three constant stages.
\end{proof}

\begin{lemma}[Strong typing survives a block of definitional equations]
\label{thm:adddefeqs-strong}
\contrib{iota}
\lean{Lean4Lean.VEnv.addDefEqs_strong}
\leanok
\uses{def:envstrong, def:ontypes, def:vdecl-wf}
\srcloc{Lean4Lean/Theory/Typing/EnvLemmas.lean}{162}
\thesisref{no thesis counterpart}

\lead{In short} \texttt{addDefEqs} \hl{leaves every constant and axiom strongly typed}, given
the start and new equations already are.
\end{lemma}
\begin{proof}
\uses{thm:ontypes-adddefeq, thm:envstrong-mono}
\leanok
\lead{Idea} Induction applying \texttt{OnTypes.addDefEq}; no \texttt{PatsStrong} hypothesis
needed, an axiom re-types nothing.
\end{proof}

\begin{lemma}[Registering $\iota$ rules preserves strong typing]
\label{thm:addrules-strong}
\contrib{iota}
\lean{Lean4Lean.VEnv.addRules_strong}
\leanok
\uses{def:envstrong, def:ontypes, family:ind-add-stages, def:ind-add-rec-rule}
\srcloc{Lean4Lean/Theory/Typing/EnvLemmas.lean}{175}
\thesisref{no thesis counterpart}

\lead{In short} \texttt{addInduct}'s rule-registration stage \hl{adds no typing obligation}.
\end{lemma}
\begin{proof}
\uses{family:indlem-foldlm, thm:ontypes-addpat, def:ind-add-rec-rule}
\leanok
\lead{Idea} Two \texttt{foldlM\_inv} reduce it to one \texttt{addRecRule}/\texttt{addPat}
step, closed by \texttt{OnTypes.addPat}; cost deferred to \ref{def:patsstrong}.
\end{proof}

\begin{lemma}[Adding the quotient block preserves strong typing]
\label{thm:addquot-strong}
\contrib{iota}
\lean{Lean4Lean.VEnv.addQuot_strong}
\leanok
\uses{def:envstrong, def:ontypes, def:venv-addquot, def:patsstrongon}
\srcloc{Lean4Lean/Theory/Typing/EnvLemmas.lean}{188}
\thesisref{axioms.tex \S2.7.1}

\lead{In short} \texttt{addQuot} \hl{preserves strong typing}, given $E$ ordered and
strongly typed, \texttt{QuotReady}, and subject reduction on the intermediate stages.
\begin{itemize}
\item \lead{Caveat} Hand-unrolled into $\sim$25 lines of \texttt{.trans} chains, since
  \texttt{addQuot} is a bind chain, not a fold, so cannot reuse
  \ref{thm:foldlm-addconst-strong}.
\end{itemize}
\end{lemma}
\begin{proof}
\uses{thm:addquot-chain, thm:ontypes-addconst, thm:ontypes-adddefeq, thm:envstrong-of-hastype, family:indlem-pats-defeqs-preserved}
\leanok
\lead{Idea} Decompose via \ref{thm:addquot-chain}, apply \texttt{OnTypes.addConst} four
times and \texttt{OnTypes.addDefEq} once; strengthen \texttt{quotDefEq} via
\texttt{EnvStrong.of\_hasType}.
\end{proof}

\begin{theorem}[Adding an inductive block preserves strong typing]
\label{thm:addinduct-strong}
\contrib{iota}
\lean{Lean4Lean.VEnv.addInduct_strong}
\leanok
\uses{def:envstrong, def:ontypes, def:ind-add-induct, structure:ind-wf, def:patsstrongon}
\srcloc{Lean4Lean/Theory/Typing/EnvLemmas.lean}{223}
\thesisref{axioms.tex \S2.6 (inductive types)}

\lead{In short} \texttt{addInduct} \hl{preserves strong typing}, given $E$ ordered and
strongly typed, \texttt{decl.WF E}, and subject reduction on its four stages.
\end{theorem}
\begin{proof}
\uses{thm:indlem-add-induct-stages, thm:foldlm-addconst-strong, thm:addrules-strong, family:indlem-stage-ordered, family:indlem-stage-le, family:indlem-pats-defeqs-preserved, structure:ind-wf}
\leanok
\lead{Idea} Split into four stages; the three constant stages go through
\ref{thm:foldlm-addconst-strong}, the rule stage through \ref{thm:addrules-strong},
re-aiming \texttt{PatsStrong} each time.
\end{proof}

\begin{theorem}[Constants and axioms of a well-formed environment are strongly typed]
\label{thm:wf-strong}
\contrib{iota}
\lean{Lean4Lean.VEnv.WF.strong}
\leanok
\uses{def:venv-wf, def:patsstrong, def:envstrong, def:ontypes}
\srcloc{Lean4Lean/Theory/Typing/EnvLemmas.lean}{263}
\thesisref{no direct counterpart; the machinery behind regularity for the annotated judgment}

\lead{In short} \texttt{E.WF} with \texttt{E.PatsStrong} gives \hl{\texttt{OnTypes E
(EnvStrong E)}}: every constant and axiom typed in the strong judgment.
\end{theorem}
\begin{proof}
\uses{def:wfprefix, thm:wfprefix-trans, thm:vdecl-wf-le, thm:ontypes-addconst, thm:ontypes-adddefeq, thm:envstrong-of-hastype, thm:foldlm-addconst-strong, thm:adddefeqs-strong, thm:addquot-strong, thm:addinduct-strong, thm:wf-ordered}
\leanok
\lead{Idea} Induction on \texttt{VEnv.WF'} generalised over a \texttt{WFPrefix} witness, so
\texttt{PatsStrong} aims at the current stage; the seven forms dispatch to
\ref{thm:foldlm-addconst-strong}, \ref{thm:adddefeqs-strong}, \ref{thm:addquot-strong},
\ref{thm:addinduct-strong}. Replaces upstream's \texttt{Ordered.strong}.
\end{proof}

\begin{theorem}[Regularity of $\iota$ reductions (open)]
\label{thm:wf-patsstrong}
\contrib{iota}
\lean{Lean4Lean.VEnv.WF.patsStrong}
\leanok
\uses{def:venv-wf, def:patsstrong}
\srcloc{Lean4Lean/Theory/Typing/EnvLemmas.lean}{334}
\thesisref{unique.tex ``Regularity of reductions'' item 1 and typesys.tex ``Regularity continued'' item red\_equiv, which the thesis dismisses as ``easy inductions''}

\lead{In short} Every well-formed environment satisfies \texttt{PatsStrong}: a well-typed
redex of a registered $\iota$ rule reduces to a \hl{reduct of the same type}.
\begin{itemize}
\item Not a consequence of \texttt{Ordered} alone: \texttt{defeq} admits any well-typed
  axiom, e.g. $\mathrm{List}\,\mathbb{N}\equiv\mathrm{List}\,\mathrm{Bool}$ gives
  \texttt{List.rec} a well-typed redex with an ill-typed reduct.
\item \lead{Caveat} The only new \texttt{sorry} on this branch: via
  \ref{thm:wf-orderedstrong}'s \texttt{CoeOut}, all 343 \texttt{VEnv.WF}-dependent constants
  -- the widest reach of any gap in the project -- become conditional on it. Upstream's entry
  point was proved.
\end{itemize}
\axfoot{sorryAx (this declaration is itself sorry), Quot.sound, propext}
\end{theorem}
\begin{proof}
\uses{thm:injectivity-open, thm:foralle-inv-derived, structure:ind-wf}
\stSorry{} \lead{Missing} Inversion of the redex's typing with injectivity of the block's
type formers (\ref{thm:injectivity-open}, itself open); a proof must also handle $\iota$ in
the presence of $\delta$ and the quotient rule.
\end{proof}

\begin{theorem}[A well-formed environment satisfies the strong-system hypotheses]
\label{thm:wf-orderedstrong}
\contrib{iota}
\lean{Lean4Lean.VEnv.WF.orderedStrong, Lean4Lean.instCoeOutWFOrderedStrong}
\leanok
\uses{def:venv-wf, struct:orderedstrong}
\srcloc{Lean4Lean/Theory/Typing/EnvLemmas.lean}{339}
\thesisref{no thesis counterpart}

\lead{In short} \texttt{E.WF} implies \hl{\texttt{E.OrderedStrong}} (ordered, strongly
typed, subject reduction for its rules); a \texttt{CoeOut} instance makes it silent.
\begin{itemize}
\item \lead{Caveat} Delivery mechanism for \ref{thm:wf-patsstrong}'s gap
  (\srcloc{Lean4Lean/Theory/Typing/EnvLemmas.lean}{342}): any \texttt{E.WF} hypothesis
  inherits the \texttt{sorry} with no visible cue.
\end{itemize}
\axfoot{sorryAx via VEnv.WF.patsStrong; Classical.choice, Quot.sound, propext}
\end{theorem}
\begin{proof}
\uses{thm:wf-ordered, thm:wf-strong, thm:wf-patsstrong, struct:orderedstrong}
\leanok
\lead{Idea} Bundle \ref{thm:wf-ordered} and \ref{thm:wf-strong} applied to
\ref{thm:wf-patsstrong}, itself at the identity prefix; only \ref{thm:wf-patsstrong} is
missing.
\end{proof}

\section{Definitional inversion}

Two long-open principles, needed by \ref{thm:wf-patsstrong}.

\begin{lemma}[Definitional inversion (open)]
\label{thm:injectivity-open}
\lean{Lean4Lean.VEnv.IsDefEqU.sort_inv, Lean4Lean.VEnv.IsDefEqU.forallE_inv_stratified, Lean4Lean.VEnv.IsDefEqU.sort_forallE_inv}
\leanok
\uses{def:hastype-istype, def:venv-wf, def:hastypestratified, def:onctx}
\srcloc{Lean4Lean/Theory/Typing/Injectivity.lean}{11}
\thesisref{unique.tex Theorem (Definitional inversion) thm:1dinv items 1--3; soundness.tex ``Type injectivity''}

\lead{In short} Over a well-formed environment and context: \hl{sorts and $\Pi$-types are
injective} under definitional equality.
\begin{itemize}
\item Equal sorts have equivalent levels; equal $\Pi$-types have equal domains/codomains
  (for the stratified \texttt{HasTypeStratified}); a sort is never equal to a $\Pi$-type.
\end{itemize}
\axfoot{sorryAx (all three are sorry), Quot.sound, propext}
\end{lemma}
\begin{proof}
\stSorry{} \lead{Missing} All three proofs; pre-existing upstream holes, not contributed, and
the ingredient \ref{thm:wf-patsstrong} names as missing.
\end{proof}

\begin{lemma}[$\Pi$-injectivity for the unstratified judgment]
\label{thm:foralle-inv-derived}
\lean{Lean4Lean.VEnv.IsDefEqU.forallE_inv}
\leanok
\uses{def:hastype-istype, def:venv-wf, def:onctx}
\srcloc{Lean4Lean/Theory/Typing/Injectivity.lean}{23}
\thesisref{unique.tex thm:1dinv item 2}

\lead{In short} In a well-formed environment with a well-formed context, equal $\Pi$-types
$\forall x{:}A.B$, $\forall x{:}A'.B'$ have \hl{equal domains and codomains}: $A\equiv A'$
and $B\equiv B'$ at some sort in $A::\Gamma$.
\axfoot{sorryAx via IsDefEqU.forallE\_inv\_stratified and via IsDefEq.strong (VEnv.WF.orderedStrong, VEnv.WF.patsStrong); Classical.choice, Quot.sound, propext}
\end{lemma}
\begin{proof}
\uses{thm:injectivity-open, thm:isdefeq-strong, def:hastypestratified}
\leanok
\lead{Idea} Strengthen via \texttt{IsDefEq.strong}, stratify, apply the admitted
\texttt{forallE\_inv\_stratified}. Byte-identical to upstream; \texttt{IsDefEq.strong}'s
environment hypothesis moved to \texttt{OrderedStrong}, so this proof now carries a second,
independent gap.
\end{proof}

\section*{Caveats}

\begin{itemize}
\item \texttt{VEnv.WF.patsStrong} (\ref{thm:wf-patsstrong}) is \texttt{sorry}: every
  \texttt{VEnv.WF} hypothesis silently inherits it via the \texttt{CoeOut} to
  \texttt{OrderedStrong} (\ref{thm:wf-orderedstrong}).
\item \ref{rule:isdefeq-pat} builds subject reduction for $\iota$ into the judgment itself;
  false over a merely \texttt{Ordered} environment.
\item The three principles of \ref{thm:injectivity-open} are \texttt{sorry} (pre-existing);
  a proof of \ref{thm:wf-patsstrong} needs them.
\item \texttt{PatTyped} (\ref{def:pattyped}) is existential and constrains only the holes a
  reduct uses: a rule can be well-formed at one convenient instantiation.
\item The \texttt{Check}/\texttt{Realizes} side-condition machinery is exercised only at the
  trivial check; no registered rule yet uses a nontrivial one.
\item \texttt{addQuot\_strong} (\ref{thm:addquot-strong}) hand-unrolls five steps rather than
  reusing \ref{thm:foldlm-addconst-strong}, since \texttt{addQuot} is a bind chain, not a
  fold.
\end{itemize}
